\chapter{Sizing of an Aerospace Helicopter Main Gearbox}
\section{Helicopter Transmission Architectures and the UH-60 Black Hawk MGB}
    In rotorcraft engineering, the Main Gearbox (MGB) is the mechanical heart of the aircraft. 
    It must safely transmit thousands of kilowatts of power from high-speed gas turbine engines 
    (typically running between 6,000 and 21,000 RPM) down to the slow, extremely high-torque 
    operational speeds dictated by the main rotor blades (typically 250 to 400 RPM). To achieve 
    this immense reduction while remaining light enough for flight, helicopter transmissions 
    universally employ a multi-stage architecture.

    The standard, most reliable MGB configuration consists of a horizontal double-input 
    engine stage transitioning through a 90° spiral bevel reduction into a vertical, high-torque 
    final planetary (epicyclic) stage as shown in \autoref{fig:gearbox6}.
    	\begin{figure}[h]
			\centering
			\includegraphics[width=0.9\textwidth]{gearbox/gearbox6}
			\caption{Helicopter MGB Configuration}
			\label{fig:gearbox6}
		\end{figure}
		\FloatBarrier

    A prime operational example of this configuration is found in the Sikorsky UH-60 Black Hawk. 
    The Black Hawk uses a modular, triple-stage main transmission. Power enters horizontally from 
    two General Electric T700 turboshaft engines into two independent input modules. Each module 
    uses high-speed helical gears to combine the engine inputs, which then drive a Spiral Bevel 
    Gear mesh. This bevel mesh executes a 90-degree axis flip, turning the horizontal engine power 
    into a vertical spinning axis. This vertical core shaft feeds directly into a final, 
    single-stage planetary gear system located beneath the main rotor mast. The planetary stage 
    uses a fixed outer ring gear and five power-splitting planet gears mounted to a forged steel 
    planet carrier, which directly rotates the main rotor shaft to lift the aircraft.

\section{Sizing Input Parameters}
    To build an MGB, every baseline physical input, material limit, and structural factor must 
    be explicitly stated. The following section defines the variables required, along with their 
    standard engineering ranges for aerospace and heavy-duty drivetrain applications:
    	\begin{enumerate}
			\item \textbf{Operational Inputs}\par
                \begin{itemize}
                    \item \textbf{P [W] (Nominal Input Power):} The total rated mechanical power delivered by the driving gas turbine engine modules.
                    \item \textbf{$n_{in}$ [RPM] (Input Engine Speed):} The rotational speed entering the horizontal input module. Typical turboshaft outputs range from 6,000 to 21,000 $RPM$.
                    \item \textbf{$n_{out}$ [RPM] (Target Main Rotor Speed):} The final required speed of the rotor mast, strictly bounded by blade radius aerodynamics to prevent supersonic tip shockwaves. Typical limits range from 250 to 400 $RPM$.
                    \item \textbf{$N_{p}$ (Number of Planet Gears):} The integer count of power-splitting planet channels evenly spaced inside the ring gear. Typical configurations utilize 3 to 6 planets; the UH-60 Black Hawk utilizes 5 planets.
                \end{itemize}

            \item \textbf{Material Strength and Density Limits}\par
                \begin{itemize}
                    \item \textbf{$\sigma_{allow,F}$ [Pa] (Allowable Bending Stress Limit):} The maximum tensile fatigue stress the tooth root can handle over millions of cycles before fracturing. For aerospace-grade carburized steels, this ranges from 400 MPa to 550 MPa.
                    \item \textbf{$\sigma_{allow,H}$ [Pa] (Allowable Contact Stress Limit):} The maximum baseline surface pressure the tooth faces can handle before pitting. For case-hardened aerospace steels, this ranges from 1,400 MPa to 1,750 MPa.
                    \item \textbf{$\sigma_{allow,shaft}$ [Pa] (Allowable Shaft Fatigue Limit):} The fully modified endurance limit of the shaft metal under reversed loading. Typical bounds range from 100 MPa to 180 MPa.
                    \item \textbf{$Z_{E}$ [\text{[Pa$^{0.5}$]}] (Material Elastic Coefficient):} Combines Young's Modulus and Poisson's ratio to govern micro-contact strain flattening.
                    \item \textbf{$\rho_{g}$ [$\text{kg}/\text{m}^3$] (Gear Material Density):} High-strength carburized alloy gear steel matches a hard constant of approximately 7,850 $\text{kg}/\text{m}^3$.
                    \item \textbf{$\rho_{s}$ [$\text{kg}/\text{m}^3$] (Shaft Material Density):} Forged steel structural shafts match a constant of approximately 7,850 $\text{kg}/\text{m}^3$.
                    \item \textbf{$\rho_{h}$ [$\text{kg}/\text{m}^3$] (Housing Material Density):} To minimize flight mass, helicopter casings utilize ultra-lightweight magnesium alloys 1,800 $\text{kg}/\text{m}^3$ or high-grade aluminum alloys 2,700 $\text{kg}/\text{m}^3$.
                \end{itemize}
            
            \item \textbf{Geometric and Structural Factors}\par
            	\begin{itemize}
                    \item \textbf{$\alpha_n$ [Degrees] (Normal Pressure Angle):} The base angle cut of the gear profile teeth. Standard industrial hobs default to ($20^\circ\text{C}$ or $25^\circ\text{C}$).
                    \item \textbf{$\beta$ [Degrees] (Helix Angle):} The angle of the helical teeth cut relative to the shaft axis. Standard industrial limits range from ($10^\circ\text{C}$ to $30^\circ\text{C}$); a setting of ($0^\circ\text{C}$) reduces the algorithm to a basic spur profile.
			        \item \textbf{$\lambda_d$ (Planetary Face-Width-to-Diameter Ratio):} The chosen aspect ratio of the gear width relative to its sun diameter. To maintain structural stiffness against shaft deflection, standard practice enforces $0.5-2.0$.
			        \item \textbf{$\lambda_b$ (Bevel Face Width-to-Diameter Ratio):} Aspect ratio bounding the spiral bevel stage cone width. Standard Aerospace Range: 0.25 to 0.35.
			        \item \textbf{$I_b$ and $I_p$ (AGMA Geometry Factors):} Dimensionless indices mapping localized involute curvature profiles for the planetary and bevel teeth respectively. Planetary Range: Evaluated via ratio kinematics (0.11 to 0.15). Spiral Bevel Range: Catalog/standard fixed constant (0.07 to 0.085).
			        \item \textbf{$C_w$ (Weight Correction Factor):} The percentage multiplier modelling material precision-machined out of the gear blank core to create lightweight flight webs. A solid blank is 1.0, while pocketed aerospace blanks range from 0.4 to 0.7.
			        \item \textbf{$C_t$ (Ring Wall Thickness Factor):} The ratio of the ring gear outer wall thickness relative to the standard module size, sized to prevent radial hoop fatigue. Valid boundaries enforce 3.5 to 6.0.
			        \item \textbf{$C_c$ (Carrier Plate Width Factor):} Sizes the width of the structural plates flanking the gears as a percentage of the face width. Standard packaging balances between 0.2 and 0.4.
			        \item \textbf{$K_a$ (Application/Overload Factor):} A designer-selected safety cushion multiplier that increases sizing limits to handle sudden engine torque surges or aerodynamic rotor blade shocks. Standard boundaries dictate1.2 to 1.5.
			        \item \textbf{$K_t$ (Temperature Factor):} An environmental fatigue scaling modifier. Set to a constant of 1.0 for all operational gearboxes running under a thermal ceiling of ($120^\circ\text{C}$).
			        \item \textbf{$n_s$ (Shaft Factor of Safety):} The explicit safety margin index applied to the shaft diameters to handle multi-plane loading uncertainty. Standard practice ranges from 1.5 to 3.0.
                \end{itemize}

            \item \textbf{Geometric and Structural Factors}\par
            	\begin{itemize}
			        \item \textbf{$t_{case}$ [m] (Perimeter Wall Thickness):} The casting wall thickness of the main outer gearbox body barrel shell. Standard bounds measure between 6 mm to 14 mm.
			        \item \textbf{$t_{plate}$ [m] (Face Plate Thickness):} The structural thickness of the casing end caps enclosing the assembly. Standard bounds range between 10 mm and 20 mm.
			        \item \textbf{$L_{clearance}$ [m] (Assembly Clearance Space):} The baseline physical gap length allowed on each end of the gear stacks to clear supporting oil seals, retaining snap-rings, and bearing inner caps. Standard spacing defaults between 20 mm and 50 mm.
                \end{itemize}
		\end{enumerate}

\section{Kinematic Speed Mapping}
    To find the weight and size of the MGB, the entire transmission is modelled as a unified 
    structural closed loop. Sizing begins by tracing how torque and rotational velocities scale 
    across the system zones based on nominal input power (P), input engine speed ($n_{in}$), bevel 
    reduction ratio ($u_{b}$), and planetary reduction ratio ($u_{p}$). 

    The total kinematic reduction ratio ($u_{overall}$) is the strict product of both internal stages:
        \begin{equation}
            \label{eq:totalkinematicreductionratio}
            u_{\text{overall}}
            \;=\;
            u_{\text{b}} \cdot u_{\text{p}}
            \;=\;
            \frac{n_{\text{in}}}{n_{\text{rotor}}}
        \end{equation}

    The intermediate velocity on the vertical transition shaft between the bevel output and planetary input is 
    calculated automatically once the optimal bevel ratio is resolved by the global optimization loop
        \begin{equation}
            \label{eq:optimalbevelratio}
            u_{\text{vertical}}
            \;=\;
            \frac{n_{\text{in}}}{n_{\text{b}}}
        \end{equation}

    The nominal torque profiles at each independent checkpoint are established by applying the AGMA 
    application overload factor ($K_{a}$) and temperature factor ($K_{t}$):
        \begin{equation}
            \label{eq:torquedesignbevela}
            T_{\text{design bevel}}
            \;=\;
            \frac{30\,P}{\pi\, n_{\text{in}}}
            \;\cdot\;
            K_{\text{a}}\;\cdot\;
            K_{\text{t}}
        \end{equation}    
 
        \begin{equation}
            \label{eq:torquedesignbevelb}
            T_{\text{design planet}}
            \;=\;
            \frac{30\,P}{\pi\, n_{\text{in}}}
            \;\cdot\;
            K_{\text{a}}\;\cdot\;
            K_{\text{t}}\;\cdot\;
            \eta_{\text{b}}\;\cdot\;
            u_{\text{b}}
        \end{equation}
    
    Where $\eta_{b}$ represents the localized mechanical transmission efficiency of the spiral bevel mesh.

\section{Mathematical Derivation of Component Weights}
    To avoid traditional \text{'guess-and-check'} verification routines, we invert the standard governing stress 
    and strain inequalities. By isolating geometric parameters (such as the standard module ($m_{std}$) and 
    pitch diameters (d)) as dependent functions of material strength limits, we can calculate physical 
    component mass directly from upfront preliminary specifications. Sizing calculations for every primary 
    component flow step-by-step from the fundamental mass-volume relationship:
        \begin{equation}
            \label{eq:massvolume}
            W
            \;=\;
            \rho\;\cdot\;
            V
        \end{equation}

    \subsection{Gears Mass}
        The total rotating gear mass is the sum of the material volumes of the spiral bevel pinion and gear, 
        plus the planetary sun, planets, and outer ring gear shells multiplied by the gear steel density $\rho_{g}$:
            \begin{equation}
                \label{eq:steeldensity}
                W_{\text{gears}}
                \;=\;
                W_{\text{bevel}} + W_{\text{planetary gears}}
            \end{equation}
    
            \begin{itemize}
                \item \textbf{Spiral Bevel Stage Sizing:}
                Spiral bevel gears are modelled as conical disks. According to standard AGMA 2003 criteria, their outer 
                pitch diameters must scale to satisfy surface contact fatigue (Hertzian pitting) thresholds under heavy 
                combined thrust loads. The minimum pinion outer pitch diameter:
                    \begin{equation}
                        \label{eq:minpinionouterpitchdia}
                        d_{\text{b,\,pinion}}
                        \;=\;
                        \left(\frac{\,2\,\;\cdot\;\;T_{\text{design bevel}}}
                        {\lambda_{\text{b}}\;\cdot\; I_{\text{b}}\;\cdot\;
                        \left(\dfrac{\sigma_{\text{allow, H, bevel}}}{Z_{\text{E}}}\right)^{\!2}}\right)^{\!1/3}
                    \end{equation}

                Integrating the conical volumes yields the combined mass envelope:
                    \begin{equation}
                        \label{eq:massenvelopea}
                        W_{\text{bevel}}
                        \;=\;
                        \bigl(\;\frac{\pi}{6}\cdot \rho_{\text{g}}\cdot C_{\text{w}}\cdot \lambda_{\text{b}}
                        \bigr)\;\cdot\;
                        \bigl(1 + u_{\text{b}}^{\,3}\bigr)\;\cdot\;
                        \bigl(d_{\text{b,\,pinion}}\bigr)^{3}
                    \end{equation}

                Substituting the outer pitch diameter boundary profile directly into the volume calculation eliminates 
                the face width ratio ($\lambda_{\text{b}}$), isolating the weight footprint strictly to raw torque and 
                material parameters:
                    \begin{equation}
                        \label{eq:massenvelopeb}
                        W_{\text{bevel}}
                        \;=\;
                        \biggl(\frac{\pi\;\cdot\;\rho_{\text{g}}\;\cdot\;C_{\text{w}}\;\cdot\;
                        Z_{\text{E}}^{\!2}\;\cdot\;
                        \bigl(1 + u_{\text{b}}^{\,3}\bigr)}{3\;\cdot\;I_{\text{b}}\;\cdot\;
                        \Bigl(\sigma_{\text{allow, H, bevel}}\Bigr)^{\!2}}\biggr)\;\cdot\;T_{\text{design bevel}}
                    \end{equation}

                \item \textbf{Planetary Stage Sizing:}
                The planetary stage consists of 1 central Sun gear (s), $N_{\text{p}}$ orbiting Planet gears (p), 
                and 1 stationary internal Ring gear (r). The base volume of the gear blanks is defined by summing 
                the individual cylindrical envelopes:
                    \begin{equation}
                        \label{eq:volumeplanetarya}
                        V_{\text{planetary}}
                        \;=\;
                        V_{\text{sun}}
                        \;+\;
                        N_{\text{p}}   \;\cdot\; V_{\text{planets}}
                        \;+\;
                        V_{\text{ring}}
                    \end{equation}

                    \begin{equation}
                        \label{eq:volumeplanetaryb}
                        V_{\text{planetary}}
                        \;=\;
                        \pi\;\cdot\;\frac{d_{\text{s}}^{\,2}}{4}\;\cdot\;FC_{\text{w}}
                        \;+\;
                        N_{\text{p}}\;\cdot\;\Bigl(\pi\;\cdot\;\frac{d_{\text{p}}^{\,2}}{4}\Bigr)\;\cdot\;FC_{\text{w}}
                        \;+\;
                        \frac{\pi}{4}\;\cdot\;\bigl(d_{\text{outer, ring}}^{\,2}\;-\;d_{\text{r}}^{\,2}\bigr)\;\cdot\;FC_{\text{w}}
                    \end{equation}

                Where F is the face width. To prevent severe radial strain from opening up the gear meshes as the planets pass by, 
                the outer ring gear envelope is modelled as a realistic hollow shell with a protective wall thickness proportional 
                to the tooth module:
                    \begin{equation}
                        \label{eq:outerringdia}
                        d_{\text{outer, ring}}
                        \;=\;
                        d_{\text{r}}
                        \;+\;
                        2C_{\text{t}}\;\cdot\;m_{\text{std}}
                    \end{equation}

                Expressing all diameters using the metric standard module (\( d = m_{\text{std}}\!\cdot\! z \))
                and substituting the coaxial kinematic rules (\( z_{\text{r}} = (u_{\text{p}}-1)\,z_{\text{s}}\))
                and (\( z_{\text{p}} = \frac{u_{\text{p}}-2}{2}\,z_{\text{s}} \)), for fixed ring gear 
                configuration, the total planetary gear weight simplifies to:
                    \begin{equation}
                        \label{eq:totalplanetarygearsmass1}
                        W_{\text{planetary gears}}
                        = \bigl(\tfrac{\pi}{4}\;\cdot \rho_{\text{g}}\;\cdot C_{\text{w}}\;\cdot \lambda_{\text{d}}\bigr)\;
                        m_{\text{std}}^{\,3}\;\cdot
                        \bigl(z_{\text{s}}^{\,3}\;
                        \bigl( 1 + N_{\text{p}}\,
                        \bigl(\tfrac{u_{\text{p}}-2}{2}\bigr)^{\,2}\bigr)
                        \;+\;
                        \bigl(4C_{\text{t}}\,(u_{\text{p}}-1)\,z_{\text{s}}^{\,2} + 4C_{\text{t}}^{\,2}\,z_{\text{s}}\bigr)\bigr)
                    \end{equation}
                
                To establish the boundaries, the required module must satisfy both AGMA root bending fatigue and contact 
                pressure simultaneously:
                    \begin{equation}
                        \label{eq:mbending}
                        m_{\text{bending}}
                        \;=\;
                        \Bigl(
                            \frac{2\,\;\cdot\;T_{\text{design planet}}}{N_{\text{p}}
                            \;\cdot\;
                            z_{\text{s}}^{\,2}
                            \;\cdot\;
                            \lambda_{\text{d}}
                            \;\cdot\;
                            \sigma_{\text{allow, F}}
                            \;\cdot\;
                            Y}
                        \Bigr)^{\,1/3}
                    \end{equation}

                    \begin{equation}
                        \label{eq:mcontact}
                        m_{\text{contact}}
                        \;=\;
                        \Bigl(
                        \frac{2\,\;\cdot\;z_{\text{E}}^{\,2}\;\cdot\;T_{\text{design planet}}}{N_{\text{p}}
                        \;\cdot\;
                        z_{\text{s}}^{\,3}
                        \;\cdot\;
                        \lambda_{\text{d}}
                        \;\cdot\;
                        I_{\text{p}}
                        \;\cdot\;
                        \sigma_{\text{allow, H, planet}}^{\,2}}
                        \Bigr)^{\,1/3}
                    \end{equation}

                Where Y is the Lewis form factor, and $\lambda_{\text{d}}$ is the face-width-to-diameter ratio (F / $d_{\text{s}}$).  

                To ensure complete non-linear closure, the sun teeth parameter is locked to its absolute geometric floor 
                to eliminate manufacturing tool undercutting:
                    \begin{equation}
                        \label{eq:zs}
                        z_{\text{s}}
                        \;=\;
                        \frac{2\,\cos^{\,3}\beta}{\sin^{\,2}\alpha_{n}}
                    \end{equation}

                The constraint is determined via the maximum envelope function \autoref{eq:maxenvelopefunction}
                which is rounded up to the nearest discrete tool cutter size from 
                the standard ISO-metric catalog \((m_{\text{std}} \ge m_{\text{calc}})\).
                    \begin{equation}
                        \label{eq:maxenvelopefunction}
                        m_{\text{calc}}
                        \;=\;
                        \max\!\bigl(m_{\text{bending}},\,m_{\text{contact}}\bigr)
                    \end{equation}

                Substituting the dominant contact-driven module expression into the cubic core completely eliminates 
                the tooth count placeholders from the baseline diameter, yielding the substituted gear mass function:
                    \begin{equation}
                    \label{eq:planetarygearsmass}
                    W_{\text{planetary gears}}
                    =
                    \left(\frac{\pi\,\rho_{\text{g}}\,C_{\text{w}}\,T_{\text{design planet}}}
                    {2\,N_{\text{p}}\,I_{\text{p}}\,
                    \sigma_{\text{allow, H, planet}}^{\,2}}\right)
                    \;\cdot
                    \bigl(\left( 1 + N_{\text{p}}
                    \left(\frac{u_{\text{p}}-2}{2}\right)^{\,2}\right)
                    + \frac{4C_{\text{t}}\,(u_{\text{p}}-1)}{z_{\text{s}}}
                    + \frac{4C_{\text{t}}^{\,2}}{z_{\text{s}}^{\,2}}\bigr)
                    \end{equation}
                \end{itemize}

    \subsection{Structural Transmission Shafts and Pins Mass}
        The MGB shaft assembly weight tracks three active loading paths: the horizontal engine input shafts, the 
        vertical transition sun shaft, and the planet pins.

        Because the planet gears encircle the central sun axis symmetrically, their radial separation vectors push 
        inward from opposite directions simultaneously, causing the net global bending moments on the main line 
        components to cancel completely to zero. Thus, the sun and carrier shafts experience pure torsion (T).

        Sized using the ASME transmission shaft design code, their safe diameters match the distortion 
        energy (von Mises) fatigue criteria: 
            \begin{equation}
                \label{eq:torsiondia}
                d_{\text{torsion}}
                \;=\;
                \biggl( \frac{\,16\,\sqrt{3}\;\cdot\;n_{\text{s}}\;\cdot\;T}
                {\pi\,\cdot\,\sigma_{\text{allow shaft}}}\biggr)^{1/3}
            \end{equation}

        The physical lengths of your shafts are strictly bound by the width of the gears mounted on them plus 
        a fixed assembly clearance ($L_{\text{clearance}}$) for bearings and oil seals: 
            \begin{itemize}
				\item Sun Shaft Length ($L_{\text{s}}$), must stretch across the gear face width:
                    \begin{equation}
                        \label{eq:sunshaftlength}
                        L_{\text{s}}
                        \;=\;
                        F
                        \;=\;
                        \lambda_{\text{d}}\;\cdot\;m_{\text{std}}\;\cdot\;z_{\text{s}}
                    \end{equation}

                \item Carrier Shaft Length ($L_{\text{c}}$), must clear the gear face width plus the thickness 
                of the two structural carrier plates capping it:
                    \begin{equation}
                        \label{eq:carriershaftlength}
                        L_{\text{c}}
                        \;=\;
                        F\;+\;2\,\cdot\,t_{\text{plate carrier}}
                        \;=\;
                        \lambda_{\text{d}}\;\cdot\;m_{\text{std}}\;\cdot\;z_{\text{s}}\;
                        \cdot\,\bigl(1\;+\;2\,\cdot\,C_{\text{c}}\bigr)
                    \end{equation}
			\end{itemize}

        By substituting the gear-blank face width definition ($F\;=\;\lambda_{\text{d}}\;\cdot\;m_{\text{std}}\;\cdot\;z_{\text{s}}$) 
        to define the physical shaft spans, we eliminate the independent length variables. The linear scaling weights 
        for the input and output shafts are defined strictly by torque, tooth count, module, and material bounds:
            \begin{equation}
                \label{eq:sunshaftmassa}
                W_{\text{sun shaft}}
                = \frac{\pi}{4}\; \rho_{\text{s}}
                \;\cdot\; \lambda_{\text{d}}
                \;\cdot\; m_{\text{std}}
                \;\cdot\; z_{\text{s}}
                \;\cdot\;
                \biggl( \frac{\,16\,\sqrt{3}\;\cdot\;n_{\text{s}}\;\cdot\;T_{\text{design bevel}}}
                {\pi\,\cdot\,\sigma_{\text{allow shaft}}}\biggr)^{2/3}
            \end{equation}

            \begin{equation}
                \label{eq:sunshaftmassb}
                W_{\text{sun shaft}}
                = \frac{\pi}{4}\; \rho_{\text{s}}
                \;\cdot\; \lambda_{\text{d}}
                \;\cdot\; m_{\text{std}}
                \;\cdot\; z_{\text{s}}
                \cdot\,\bigl(1\;+\;2\,\cdot\,C_{\text{c}}\bigr)
                \;\cdot\;
                \biggl( \frac{\,16\,\sqrt{3}\;\cdot\;n_{\text{s}}\;\cdot\;u_{\text{p}}\;\cdot\;T_{\text{design bevel}}}
                {\pi\,\cdot\,\sigma_{\text{allow shaft}}}\biggr)^{2/3}
            \end{equation}
        
        \textit{Note: The only unknown left inside these bracket structures is the planetary stage 
        gear ratio ($u_{\text{p}}$), as the module ($m_{\text{std}}$), and sun teeth ($z_{\text{s}}$) are resolved beforehand.}

        Conversely, the individual planet pins are not balanced. Each pin acts as a short, localized beam of 
        length F supporting a heavy mechanical bending load. The net radial force concentrated on each pin is:
            \begin{equation}
                \label{eq:pinmass}
                W_{\text{pin}}
                \;=\;
                \frac{4T_{\text{design planet}}}
                {m_{\text{std}}\;\cdot\;z_{\text{s}}\;\cdot\;N_{\text{p}}}
            \end{equation}
        
        Treating the pin as supported on both ends by the rigid carrier walls, the peak internal 
        bending moment occurs at the center of the pin span ($M_{\text{pin}}\;=\frac{\;W_{\text{pin}}\;\cdot\;F}{4}$) 
        simplifies directly to:
            \begin{equation}
                \label{eq:pinmoment}
                M_{\text{pin}}
                \;=\;
                \frac{\lambda_{\text{d}}\;\cdot\;T_{\text{design planet}}}{N_{\text{p}}}
            \end{equation}

        Sized under pure bending fatigue (T = 0), the total pin mass across all ($N_{\text{p}}$) paths scales as: 
            \begin{equation}
                \label{eq:pinsmass}
                W_{\text{pins}} =
                \frac{\pi}{4}\;
                \cdot\; \rho_{\text{s}}
                \;\cdot\; N_{\text{p}}
                \;\cdot\; \lambda_{\text{d}}
                \;\cdot\; m_{\text{std}}
                \;\cdot\; z_{\text{s}}
                \;\cdot\;
                \biggl(\frac{\,32\;\cdot\;n_{\text{s}}\;\cdot\;
                \bigl(\frac{\lambda_{\text{d}}\;\cdot\;T_{\text{design planet}}}{N_{\text{p}}}\bigr)}
                {\pi\,\cdot\,\sigma_{\text{allow shaft}}}\biggr)^{2/3}
            \end{equation}

    \subsection{Rigid Planet Carrier Mass}
        The planet carrier is a structural frame flanking the gears that collects the planet pin loads to rotate 
        the main rotor mast. To protect the gear meshes from edge-loading stress concentrations, it must be highly 
        rigid against torsional twist strain. Sized as twin circular plates of diameter matching the outer 
        ring boundary ($d_{\text{outer ring}}$) and plate thicknesses proportional to 
        the gear width ($t_{\text{plate carrier}}\;=\;C_{\text{c}}\;\cdot\;F$), its structural mass is 
        derived directly from its geometric envelope:
            \begin{equation}
                \label{eq:carriermassa}
                W_{\text{carrier}}
                = 2\,\bigl(\frac{\pi}{4}\,d_{\text{outer ring}}^{2}
                \bigl(C_{\text{c}}\!\cdot\!F\bigr)\bigr)
                \;\cdot\;\rho_{\text{s}}\;\cdot\;K_{\text{web}}
            \end{equation}
    
        Where $K_{\text{web}}$ is the geometric modifier accounting for weight-reduction pockets 
        pocketed out of the carrier web plates \(0.60 \leq K_{\text{web}} \leq 0.75\).
        
        By substituting the gear-blank face width and pitch diameter parameters under the assumption 
        that contact stress governs the active sizing envelope, the carrier mass equation simplifies to:
            \begin{equation}
                \label{eq:carriermassb}
                W_{\text{planetary gears}}
                =
                \bigl(\frac{\pi \;\cdot\; \rho_{\text{s}}
                \;\cdot\; C_{\text{c}}
                \;\cdot\; K_{\text{web}}
                \;\cdot\; Z_{\text{E}}^{2}
                \;\cdot\; {(u_{\text{p}}-1)}^{2}}
                {N_{\text{p}}\;\cdot\;I_{\text{p}}\;\cdot\;\sigma_{\text{allow, H, planet}}^{2}}
                \bigr)\;\cdot\; T_{\text{design planet}}
            \end{equation}

    \subsection{Bearings Mass}
        Sized using an analytical load-scaling law calibrated from aerospace bearing catalogs, the mass of 
        rolling elements scales with the net localized radial force vector raised to the 1.2 power. 
        Expressing the radial forces on the planet pins and the main lines directly via torque, speed, and 
        standard module variables results in the combined bearing mass expression:
            \begin{equation}
                \label{eq:bearingsmass}
                W_{\text{bearings}}
                = N_{\text{p}}\;\cdot\;
                \bigl(X_{\text{planet}}\;\cdot\;
                \bigl{(\frac{4\;\cdot\;T_{\text{design planet}}}
                {m_{\text{std}}\;\cdot\;z_{\text{s}}\;\cdot\;N_{\text{p}}}\bigr)}^{1.2}\bigr)
                + X_{\text{shaft}}\;\cdot\;
                \bigl{(u_{\text{overall}}\;\cdot\;T_{\text{design bevel}}\bigr)}^{0.8}
            \end{equation}
        
        Where $X_{\text{planet}}$ and $X_{\text{shaft}}$ represent empirical lookup constants calibrated 
        relative to the primary shaft dimensions.

    \subsection{Structural Housing Casing Mass}
        The casing forms the environmental boundary layer, structural load path, and pressure shell enclosing 
        both the spiral bevel input meshes and the final planetary stage. To minimize structural dead-weight, 
        the housing is modeled as a unified, tapered Conical Tower Shell rather than a uniform cylinder, 
        matching the tapering force profile of the internally stacked stages.

        Sized using an ultra-lightweight magnesium alloy, the total housing mass is calculated as 
        the sum of the upper planetary cylinder barrel, the lower tapered bevel cone frustum shell, and 
        the structural sealing end plates:
            \begin{equation}
                \label{eq:housingmass}
                W_{\text{housing}}
                = \rho_{\text{h}}\;\cdot\;\pi\;\cdot\;K_{\text{rib}}\;\cdot\;
                \!\!
                \begin{aligned}[t]
                    &\bigl(\bigl(D_{\text{h, planet}}\;\cdot\;t_{\text{case}}\;\cdot\;L_{\text{h, planet}}\bigr)
                    \\[4pt]
                    &+\,\bigl(t_{\text{case}}\;\cdot\;L_{\text{h,bevel}}
                    \;\cdot\;\frac{D_{\text{h, planet}}+d_{\text{b,pinion}}}{2}\bigr)
                    \\[4pt]
                    &+\,\bigl(\tfrac{t_{\text{plate}}}{2}\;\cdot\;D_{\text{h, planet}}^{\,2}\bigr)\bigr)
                \end{aligned}
            \end{equation}

        Where the active boundary diameter and length profiles are mapped strictly from your standard 
        module and integer tooth boundaries:
			\begin{itemize}
				\item Planetary Upper Housing Diameter:
                    \begin{equation}
                        \label{eq:planetaryupperhousingdia}
                        D_{\text{h, planet}}
                        = \bigl( m_{\text{std}}\!\cdot\! z_{\text{s}}
                        + 2\,C_{\text{t}}\!\cdot\! m_{\text{std}} \bigr)
                        + 2\,t_{\text{case}}
                    \end{equation}

				\item Planetary Stage Housing Length:
                    \begin{equation}
                        \label{eq:planetarystagehousinglength}
                        L_{\text{h, planet}}
                        = \lambda_{\text{d}}\cdot m_{\text{std}}\cdot z_{\text{s}}\cdot
                        \bigl(1+2\,C_{\text{c}}\bigr) + L_{\text{clearance}}
                    \end{equation}

				\item Bevel Stage Housing Length:
                    \begin{equation}
                        \label{eq:bevelstagehousinglength}
                        L_{\text{h, bevel}}
                        = d_{\text{b, pinion}} \;\cdot\; \tan\!\bigl(\delta_{\text{gear}}\bigr)
                        \;+\; 2 \;\cdot\; L_{\text{clearance}}
                    \end{equation}
			\end{itemize}
        
        The factor ($K_{\text{rib}} \approx 1.20$) is the structural multiplier that adds the 
        physical weight of the deep external grid-stiffness ribs cast onto both the cylinder and 
        conical perimeter walls to prevent casing deflection under flight loads.

\section{Global Multi-Stage Ratio Optimization Split}
    The global multi-stage ratio optimization ensures minimum system mass by deriving a 
    weight minimization penalty function, focusing on the bevel ratio ($u_{\text{b}}$) and planetary ratio ($u_{\text{p}}$). 
    By expanding the derivative of this system function with respect to $u_{\text{b}}$, a master optimization of 
    quintic polynomial is established.

    From our previous derivations, we proved that the physical mass of a gear blank scales directly 
    with its torque capacity, which is controlled by its gear ratio.
		\begin{enumerate}
            \item The Spiral Bevel Stage Weight: Scales cubically with its own ratio because 
            it is a solid conical wheel setup: $W_{\text{b}}\propto u_{\text{b}}^{3}$.
            
            \item The Planetary Stage Ring Gear Weight: The outer ring gear envelope dominates the 
            final planetary weight. Its diameter scales linearly with its ratio, and its volume 
            scales cubically with that diameter: $W_{\text{ring}}\propto {\bigl(u_{\text{p}}-1\bigr)}^{3}$.

            Because the planetary stage is downstream, it receives the torque multiplied 
            by the first stage. Therefore, we must multiply the planetary weight by the 
            input stage ratio ($u_{\text{b}}$) to account for the torque scale-up:
            $W_{\text{planetary scaled}}\propto u_{\text{b}}\;\cdot\;{\bigl(u_{\text{p}}-1\bigr)}^{3}$.
            
            \item Shaft Assembly Mass:
                \begin{itemize}
                    \item Bevel Shafts: $\propto u_{\text{b}}^{2/3}$
                    \item Planetary Shafts: Constant relative to the split
                    \item Pins: $W_{\text{pins}}\propto u_{\text{b}}^{4/3}$
                \end{itemize}

            \item Planet Carrier Mass: $\propto u_{\text{b}}\;\cdot\;{\bigl(\frac{u_{\text{overall}}}{u_{\text{b}}} - 1\bigr)}^{2}$

            \item Bearings Mass: Constant relative to the split

            \item Housing Casing Mass: $\propto \frac{u_{\text{overall}}}{u_{\text{b}}}$
        \end{enumerate}

    If we add the structural scaling metrics of the housing, carrier, and shafts to our baseline 
    gear blank equations, the comprehensive mass penalty function $f(u_{\text{b}})$ expands to:
        \begin{equation}
            \label{eq:masspenaltyfunction}
            f(u_{\text{b}})
            = \begin{aligned}[t]
                &(1 + u_{\text{b}}^{3}) \\[4pt]
                &\quad + u_{\text{b}}\;
                \Bigl(\frac{u_{\text{overall}}}{u_{\text{b}}} - 1 \Bigr)^{3} \\[4pt]
                &\quad + \gamma_{1}\,u_{\text{b}}^{\frac{2}{3}}
                + \gamma_{2}\,u_{\text{b}}\;
                \Bigl(\frac{u_{\text{overall}}}{u_{\text{b}}} - 1 \Bigr)^{2} \\[4pt]
                &\quad + \gamma_{3}\;
                \Bigl(\frac{u_{\text{overall}}}{u_{\text{b}}}\Bigr)
                + \gamma_{4}\,u_{\text{b}}^{\frac{4}{3}}
            \end{aligned}
        \end{equation}

    Where $\gamma_{1}$, $\gamma_{2}$, $\gamma_{3}$, and $\gamma_{4}$ are structural mass distribution 
    constants determined by your material choices and geometric properties.

\section{Analytical Minimization and Derivation of the Master Quintic Polynomial}
    To resolve the precise architectural configuration that yields the absolute minimum mass 
    for the entire assembled Main Gearbox (MGB), calculus must be applied to the comprehensive system function.

    To locate the absolute minimum mass floor, we evaluate the first derivative of the expanded system mass 
    function with respect to our primary design variable ($u_{\text{b}}$), setting the resulting 
    expression equal to zero: $\frac{\partial f}{\partial u_{\text{b}}}=0$.
    
    Once your optimization code calculates $u_{\text{b}}$, the engineering sequence unlocks completely.

    The software can now run down the line to calculate the exact tooth sizes, shaft dimensions, 
    and material volumes for every single component in your transmission tower. 
    A final linear summation of all five computed sub-assembly blocks to output the exact flight weight 
    of the finished transmission assembly.

\section{Comprehensive Housing Size}
    To calculate the spatial bounding box of the complete, fully enclosed helicopter main gearbox tower, 
    we translate our discrete design coordinates ($m_{\text{std}}$, $z_{\text{s}}$, $z_{\text{r}}$, and $d_{\text{b, pinion}}$) 
    straight into a physical three-dimensional envelope. Because the architecture wraps both the narrower 
    horizontal spiral bevel input and the wider vertical planetary cell into a single, cohesive structural shell, 
    the geometric package is completely defined by two primary parameters: the Maximum System 
    Outer Housing Diameter ($d_{\text{h, planet}}$) and the Total Assembly Tower Height ($H_{\text{sys}}$).

    \subsection{The Maximum System Outer Diameter}
        The maximum width across the circular perimeter profile of the transmission occurs at the 
        upper cylindrical section flanking the stationary planetary ring gear:
            \begin{equation}
                \label{eq:maxsystemouterdia}
                D_{\text{h, planet}}
                = m_{\text{std}}\!\cdot\! z_{\text{s}}\!\cdot\!
                \bigl( (u_{\text{p}}-1)
                + \frac{2\,C_{\text{t}}}{z_{\text{s}}} \bigr)
                + 2\,t_{\text{case}}
            \end{equation}

    \subsection{The Total Housing Tower Height}
        The total axial length of the gearbox shell represents the direct stacking of the vertical 
        planetary stage thickness, the spiral bevel cone drop length, and the required internal 
        bearing arrangement clearances.
            \begin{equation}
                \label{eq:totalhousingtowerheighta}
                H_{\text{sys}}
                = L_{\text{h, planet}} + L_{\text{h, bevel}}
            \end{equation}

            \begin{equation}
                \label{eq:totalhousingtowerheightb}
                H_{\text{sys}}
                = \bigl(\lambda_{\text{d}}\;\cdot\; m_{\text{std}}\;\cdot\; z_{\text{s}}
                \;\cdot\; \bigl(1 + 2\;\cdot\;C_{\text{c}}\bigr)
                +  L_{\text{clearance}}\bigr)
                + \bigl(d_{\text{b, pinion}}\;\cdot\;\tan\delta_{\text{gear}}
                + 2\;\cdot\; L_{\text{clearance}}\bigr)
            \end{equation}

\newpage
    \subsection{Total Volumetric Size}
        By treating the casing structure as an upper cylinder intersecting a lower tapered cone 
        frustum shell, the literal volumetric space occupied is:
            \begin{equation}
            \label{eq:totalvolumetricsize}
            V_{\text{envelope}}
            =
            \frac{\pi}{4}\;\cdot\;D_{\text{h, planet}}^{2}\;\cdot\;L_{\text{h, planet}}
            \;+\;
            \frac{\pi}{12}\;\cdot\;L_{\text{h, bevel}}\;\cdot\;
                \bigl(
                    D_{\text{h, planet}}^{2}
                    \;+\;
                    D_{\text{h, planet}}\;\cdot\;d_{\text{b, pinion}}
                    \;+\;
                    d_{\text{b, pinion}}^{2}
                \bigr)
            \end{equation}
